\documentclass[12pt]{article}
\usepackage[utf8]{inputenc}
\usepackage[margin=1in]{geometry}
\usepackage{amsmath}
\usepackage{graphicx}
\usepackage{booktabs}
\usepackage{hyperref}
\usepackage[round]{natbib}

\title{Modulation of Alpha Oscillations by Attention Is Predicted by Hemispheric Asymmetry of Subcortical Regions}
\author{Anonymous Authors}
\date{}

\begin{document}
\maketitle

\begin{abstract}
Spatial attention modulates posterior alpha-band (8--12~Hz) oscillations, producing contralateral power decreases and ipsilateral increases relative to the attended hemifield. While cortical contributions to this process are well characterized, the role of subcortical structures remains poorly understood, largely because magnetoencephalography (MEG) cannot reliably detect deep-source oscillatory activity. Here we used a novel indirect approach combining structural MRI volumetrics with MEG-measured alpha power to test whether hemispheric asymmetries in subcortical structures predict individual differences in alpha lateralization during spatial attention. In 33 right-handed participants who performed a cued spatial attention change-detection task, we found that the lateralization volumes of the thalamus, caudate nucleus, and globus pallidus jointly predicted hemispheric lateralization of alpha modulation (GLM $p = 0.0007$). Critically, these subcortical structures contributed differentially depending on task demands: the thalamus predicted alpha modulation under low-demand conditions, the globus pallidus at intermediate load, and the caudate nucleus under the most demanding condition involving high perceptual load with a salient distractor. These findings reveal a previously unappreciated link between subcortical structural asymmetry and functional alpha lateralization in humans, with distinct subcortical circuits engaged as task demands increase.
\end{abstract}

\section{Introduction}

Spatial attention enhances the processing of stimuli at attended locations while suppressing interference from unattended locations. A well-established neural correlate of this process is the modulation of posterior alpha-band oscillations (8--12~Hz): alpha power decreases over cortical regions contralateral to the attended hemifield and increases ipsilaterally \citep{worden2000anticipatory, thut2006alpha, klimesch2007eeg}. This hemispheric lateralization of alpha power reflects the deployment of anticipatory attention and has been shown to predict behavioral performance, with stronger lateralization associated with better target detection \citep{thut2006alpha}.

Extensive research has characterized the cortical networks responsible for controlling spatial attention, including the dorsal attention network comprising the frontal eye fields and the intraparietal sulcus \citep{corbetta2002control}. However, the contribution of subcortical structures to attention-related alpha modulation in humans remains poorly understood. Animal studies have established that subcortical nuclei, particularly the pulvinar nucleus of the thalamus, play a critical role in regulating attention and alpha oscillations \citep{saalmann2012pulvinar}. The basal ganglia, including the globus pallidus and caudate nucleus, have also been implicated in attention through their role in action selection and cognitive control \citep{desimone1995neural}. Furthermore, subcortical atrophy in structures such as the thalamus and caudate nucleus is a hallmark of neurological disorders including Parkinson's disease, Alzheimer's disease, and Huntington's disease, suggesting that oscillatory biomarkers linked to subcortical integrity could have clinical relevance.

A major methodological barrier to studying subcortical contributions is that MEG and EEG have low sensitivity to deep-source activity from the thalamus and basal ganglia due to their distance from scalp sensors and the partial cancellation of their field patterns. This limitation has left a substantial gap in our understanding of how subcortical structures shape cortical alpha dynamics in humans.

To bypass this limitation, we employed a novel indirect approach: we tested whether structural hemispheric asymmetries in subcortical volumes, measured with MRI and quantified using FreeSurfer segmentation \citep{fischl2002whole}, predict the degree of hemispheric lateralization in MEG-measured alpha power during spatial attention \citep{halassa2014thalamic}. The rationale is that if a subcortical structure contributes to alpha modulation in a lateralized fashion, then individuals with greater structural asymmetry (i.e., a larger left-than-right or right-than-left volume) should show correspondingly lateralized functional effects. We examined five bilateral subcortical structures---thalamus, caudate nucleus, putamen, globus pallidus, and nucleus accumbens---across four conditions of a factorial design that crossed target perceptual load (high vs.\ low) with distractor salience (high vs.\ low).

\section{Methods}

\subsection{Participants}

Thirty-five right-handed healthy volunteers were recruited. Two were excluded due to consent withdrawal and MRI quality issues, yielding a final sample of 33 participants (25 female, 8 male; mean age $24 \pm 5.7$ years; all with normal or corrected-to-normal vision). Participants were compensated at 15~GBP per hour. The study was approved by the STEM committee at the University of Birmingham.

\subsection{Experimental Design}

Participants performed a cued spatial attention change-detection task during MEG recording. On each trial, a central fixation cross appeared for 1000~ms, followed by bilateral face stimuli (8 identities, presented at 7 degrees eccentricity, 8 degrees diameter) for 1000~ms. A directional cue (350~ms) indicated the target hemifield. After a variable delay interval (1000--2000~ms), the eye-gaze of each face shifted (150~ms), and participants reported the target gaze direction within a 1000~ms response window.

A $2 \times 2$ factorial manipulation crossed target perceptual load (low = clear face, high = noisy face with 50\% pixel swap) with distractor salience (low = noisy distractor, high = clear distractor), producing four conditions (Table~\ref{tab:conditions}). Participants completed 2 blocks of 256 trials each (512 total) across approximately 45 minutes.

\begin{table}[ht]
\centering
\caption{Experimental conditions in the $2 \times 2$ factorial design.}
\label{tab:conditions}
\begin{tabular}{clllll}
\toprule
Condition & Target Load & Distractor Salience & Target & Distractor \\
\midrule
1 & Low & Low & Clear face & Noisy face \\
2 & High & Low & Noisy face & Noisy face \\
3 & Low & High & Clear face & Clear face \\
4 & High & High & Noisy face & Clear face \\
\bottomrule
\end{tabular}
\end{table}

\subsection{MEG Acquisition and Analysis}

MEG data were recorded using a 306-channel Elekta/MEGIN Neuromag system (204 planar gradiometers, 102 magnetometers) at a 1000~Hz sampling rate with an online bandpass filter of 0.1--330~Hz. Head position was monitored via five HPI coils, and head shape was digitized using a Polhemus Isotrak system. Data were preprocessed using spatiotemporal signal space separation \citep{taulu2006spatiotemporal}.

Alpha power modulation was quantified in the $-850$ to 0~ms pre-target interval using a modulation index:
\begin{equation}
MI(\alpha) = \frac{P_{\text{attend-left}} - P_{\text{attend-right}}}{P_{\text{attend-left}} + P_{\text{attend-right}}}
\label{eq:mi}
\end{equation}
where $P$ denotes alpha-band power. Five symmetric sensor clusters per hemisphere showing the strongest modulation were selected as regions of interest (10 sensors total).

The hemispheric lateralization modulation was computed as:
\begin{equation}
HLM(\alpha) = \frac{|MI_{\text{left}}| - |MI_{\text{right}}|}{|MI_{\text{left}}| + |MI_{\text{right}}|}
\label{eq:hlm}
\end{equation}
capturing individual differences in the relative strength of alpha modulation between hemispheres. No significant group-level hemispheric bias was observed ($p = 0.39$), indicating substantial individual variability.

\subsection{Structural MRI and Volumetrics}

Structural MRIs were acquired on a 3T scanner and segmented using FreeSurfer \citep{fischl2002whole}. Bilateral volumes were extracted for five subcortical structures: thalamus, caudate nucleus, putamen, globus pallidus, and nucleus accumbens. Lateralization volume (LV) was computed as:
\begin{equation}
LV = \frac{V_{\text{left}} - V_{\text{right}}}{V_{\text{left}} + V_{\text{right}}}
\label{eq:lv}
\end{equation}

\subsection{Statistical Analysis}

A generalized linear model (GLM) tested whether subcortical LV values predicted $HLM(\alpha)$. Model selection used the Akaike Information Criterion (AIC) and Bayesian Information Criterion (BIC) across all possible combinations of the five subcortical regressors. A multivariate regression further examined how the selected structures predicted alpha lateralization separately in each condition. Bayes factors \citep{kass1995bayes} were computed for supplementary analyses involving RIFT signals and behavioral asymmetry.

\section{Results}

\subsection{Subcortical Volume Lateralization}

Among the five subcortical structures, only the caudate nucleus showed significant volumetric lateralization, being right-lateralized ($p = 0.021$). The putamen, nucleus accumbens, and globus pallidus showed no significant lateralization. The thalamus showed a non-significant trend toward left lateralization (Table~\ref{tab:lateralization}).

\begin{table}[ht]
\centering
\caption{Subcortical volume lateralization results.}
\label{tab:lateralization}
\begin{tabular}{lccc}
\toprule
Structure & Lateralization & $p$-value & Significant \\
\midrule
Caudate nucleus & Right & 0.021 & Yes \\
Putamen & None & $> 0.05$ & No \\
Nucleus accumbens & None & $> 0.05$ & No \\
Globus pallidus & None & $> 0.05$ & No \\
Thalamus & Trend left & $> 0.05$ & No \\
\bottomrule
\end{tabular}
\end{table}

\subsection{GLM Model Selection}

AIC and BIC model comparison across all possible regressor combinations identified the three-regressor model including thalamus, caudate nucleus, and globus pallidus as the optimal model (Table~\ref{tab:modelsel}). This model significantly predicted $HLM(\alpha)$ compared to the null model ($p = 0.0007$). Single-structure models, two-structure models, and the full five-structure model all yielded higher (worse) AIC and BIC values than the selected three-structure model.

\begin{table}[ht]
\centering
\caption{GLM model selection (representative subset). The best model is highlighted.}
\label{tab:modelsel}
\begin{tabular}{lccl}
\toprule
Model Regressors & AIC & BIC & Selected \\
\midrule
Thalamus only & Higher & Higher & No \\
Caudate only & Higher & Higher & No \\
Globus pallidus only & Higher & Higher & No \\
Thalamus + Caudate & Lower & Lower & No \\
Thalamus + Globus pallidus & Lower & Lower & No \\
\textbf{Thalamus + Caudate + GP} & \textbf{Lowest} & \textbf{Lowest} & \textbf{Yes} \\
All 5 structures & Higher & Higher & No \\
\bottomrule
\end{tabular}
\end{table}

\subsection{Best Model: Pooled Results}

In the selected three-regressor model, all three subcortical structures showed significant positive beta coefficients, indicating that larger left-relative-to-right volume was associated with stronger left-hemisphere-dominant alpha modulation. The overall model was highly significant ($p = 0.0007$). Individual scatter plots confirmed positive correlations between each structure's lateralization volume and $HLM(\alpha)$.

\subsection{Condition-Specific Contributions}

Multivariate regression across the four conditions revealed a striking dissociation in the contributions of the three subcortical structures (Table~\ref{tab:conditions_results}). Under the least demanding condition (low load, low distractor salience), the thalamus was the significant predictor. Under intermediate-demand conditions (high load with low salience, or low load with high salience), the globus pallidus emerged as the dominant predictor. Under the most demanding condition (high load with salient distractor), the caudate nucleus was the significant predictor. The overall multivariate model was significant at $p = 0.001$.

\begin{table}[ht]
\centering
\caption{Subcortical contributions to alpha lateralization by task condition.}
\label{tab:conditions_results}
\begin{tabular}{llccc}
\toprule
Condition & Demand Level & Thalamus & Caudate & Globus Pallidus \\
\midrule
1: Low load, low salience & Lowest & Sig.\ (+) & n.s. & n.s. \\
2: High load, low salience & Intermediate & n.s. & n.s. & Sig.\ (+) \\
3: Low load, high salience & Intermediate & n.s. & n.s. & Sig.\ (+) \\
4: High load, high salience & Highest & n.s. & Sig.\ (+) & n.s. \\
\bottomrule
\end{tabular}
\end{table}

This dissociation suggests a hierarchical recruitment of subcortical structures as attentional demands increase. The thalamus, which receives direct sensory input, may support basic attentional filtering under easy conditions. The globus pallidus, a key output nucleus of the basal ganglia, appears to contribute when task demands require more active suppression of either perceptual noise or salient distractors. The caudate nucleus, an input nucleus of the basal ganglia involved in cognitive control, is recruited when both perceptual load and distractor salience are high, requiring the most demanding form of attentional control.

\section{Discussion}

We demonstrated that structural hemispheric asymmetries in subcortical nuclei predict functional hemispheric lateralization of attention-related alpha oscillations measured with MEG. This finding establishes, for the first time in humans, a direct link between subcortical structural organization and the well-known alpha lateralization phenomenon that accompanies spatial attention deployment \citep{worden2000anticipatory, thut2006alpha}.

The novel indirect approach employed here---correlating structural volumetric asymmetries with functional alpha lateralization---circumvents the fundamental limitation that MEG cannot reliably measure oscillatory activity from deep subcortical sources. While this approach cannot provide the same temporal or causal specificity as direct recording or stimulation of subcortical nuclei, it reveals structure-function relationships that would be invisible to either structural imaging or MEG alone.

The condition-specific dissociation among the three subcortical structures provides a nuanced view of how subcortical circuits may support attention under varying levels of demand. The thalamus, known for its role in sensory gating and attention filtering \citep{saalmann2012pulvinar, halassa2014thalamic}, dominated the easy condition where straightforward attentional allocation suffices. This is consistent with the thalamus acting as a relay that biases sensory processing in favor of the attended hemifield when the task imposes minimal additional demands.

The globus pallidus emerged as the key predictor at intermediate demand levels, consistent with its role as a primary output nucleus of the basal ganglia that tonically inhibits thalamic relay nuclei \citep{desimone1995neural}. When either the target is degraded (high load) or the distractor is salient, the globus pallidus may modulate thalamocortical circuits to enhance target processing or suppress distractor interference, adjusting the gain of alpha oscillation generators in posterior cortex.

The caudate nucleus, which dominates under the most demanding condition, is an input nucleus of the basal ganglia that receives extensive projections from prefrontal cortex and is implicated in cognitive flexibility and rule-based behavior \citep{zhong2014causal}. Its engagement in the high-load, high-salience condition suggests that when attentional demands exceed the capacity of thalamocortical and pallidal filtering mechanisms, prefrontal-caudate circuits are recruited to implement more effortful, top-down attentional control.

The right lateralization of the caudate nucleus ($p = 0.021$) is consistent with previous reports and may relate to the known right-hemisphere dominance in spatial attention \citep{corbetta2002control}. However, the globus pallidus and thalamus did not show significant group-level lateralization, suggesting that their functional contributions to alpha lateralization emerge from individual differences in volumetric asymmetry rather than from a consistent population-level structural bias.

Several limitations should be considered. First, the sample size of 33 participants, while adequate for detecting the observed effects, limits the power to detect smaller structure-function associations. Second, the correlational nature of the approach cannot establish causal relationships between subcortical volumes and alpha modulation. Third, MRI volumetrics provide only a coarse measure of subcortical organization that cannot distinguish between nuclear subdivisions. Fourth, the use of FreeSurfer automated segmentation, while highly reproducible, may introduce systematic biases in the estimation of small subcortical structures such as the nucleus accumbens.

Future work should combine this structural approach with more direct measures of subcortical function, such as high-field fMRI of thalamic and basal ganglia activity during attention tasks, or intracranial recordings in clinical populations with electrode coverage of these structures. Testing whether the observed structure-function relationships are altered in neurological conditions characterized by subcortical atrophy could establish the clinical relevance of these findings.

\section{Conclusion}

We show that hemispheric volumetric asymmetries in the thalamus, caudate nucleus, and globus pallidus predict individual differences in attention-related alpha lateralization measured with MEG. Different subcortical structures dominate under different levels of attentional demand, suggesting hierarchical recruitment of subcortical circuits as task difficulty increases. These findings reveal a previously hidden link between subcortical brain structure and the cortical oscillatory dynamics that support spatial attention, with potential implications for understanding attention deficits in subcortical neurological disorders.

\bibliographystyle{plainnat}
\bibliography{references}

\end{document}
